% Anexo A — Esquemas de cobro del impuesto predial
% Generado por scripts/build_anexo_esquemas.py
% Requiere: \usepackage{booktabs}

\subsection*{Tarifa al millar}
Aplica una tasa fija por cada millar (\textperthousand) del valor catastral, diferenciada por tipo de predio (urbano, rústico, etc.). El impuesto es lineal en el valor: $T = \tau \cdot V / 1000$.
% tarifa_millar — Tarifa al millar: San Juan del Rio, Querétaro (2015).
\begin{table}[htbp]
\centering
\small
\begin{tabular}{llrr}
\toprule
Grupo & Descripción & Tasa (al millar) & Cuota fija ad. \\
\midrule
urbano & Predio urbano edificado & 1.6 & — \\
urbano & Predio urbano baldío & 8 & — \\
rustico & Predio rústico & 1.2 & — \\
urbano & Predio de fraccionamiento en proceso de ejecución & 1.6 & — \\
urbano & Predio de reserva urbana & 1.4 & — \\
rustico & Predio de producción agrícola, con dominio pleno que provenga de ejido & 0.2 & — \\
\bottomrule
\end{tabular}
\caption{Tarifa al millar: San Juan del Rio, Querétaro (2015).}
\label{tab:esquema-tarifa_millar}

\end{table}


\subsection*{Progresivo (tasa marginal)}
Tabla por rangos de valor catastral en la que cada tramo agrega una cuota fija más una tasa marginal sobre el excedente del límite inferior: $T = c_i + \tau_i\,(V - L_i)$. La carga efectiva crece con el valor.
% progresivo — Progresivo (tasa marginal): Zapopan, Jalisco (2023).
\begin{table}[htbp]
\centering
\small
\begin{tabular}{rrrrr}
\toprule
Rango & Límite inferior & Límite superior & Cuota fija & Tasa marginal \\
\midrule
1 & \$0.01 & \$620{,}100 & \$0 & 0.00023 \\
2 & \$620{,}100 & \$944{,}864 & \$142.63 & 0.00023 \\
3 & \$944{,}864 & \$1{,}304{,}860 & \$217.33 & 0.0002495 \\
4 & \$1{,}304{,}860 & \$1{,}875{,}923 & \$307.15 & 0.0002691 \\
5 & \$1{,}875{,}923 & \$2{,}680{,}270 & \$460.83 & 0.0002886 \\
6 & \$2{,}680{,}270 & \$3{,}707{,}349 & \$692.97 & 0.0003082 \\
7 & \$3{,}707{,}349 & \$5{,}158{,}931 & \$1{,}009.52 & 0.0003277 \\
8 & \$5{,}158{,}931 & \$7{,}642{,}490 & \$1{,}485.21 & 0.0003473 \\
9 & \$7{,}642{,}490 & \$12{,}643{,}390 & \$2{,}347.76 & 0.0003668 \\
10 & \$12{,}643{,}390 & \$26{,}960{,}590 & \$4{,}182.10 & 0.0003864 \\
11 & \$26{,}960{,}590 & \$71{,}716{,}206 & \$9{,}714.27 & 0.0004059 \\
12 & \$71{,}716{,}206 & \$174{,}033{,}618 & \$27{,}880.58 & 0.0004255 \\
13 & \$174{,}033{,}618 & \$437{,}096{,}000 & \$71{,}416.64 & 0.000445 \\
14 & \$437{,}096{,}000 & \$1{,}251{,}787{,}029 & \$188{,}479.40 & 0.0004646 \\
15 & \$1{,}251{,}787{,}029 & En adelante & \$566{,}984.86 & 0.0004841 \\
\bottomrule
\end{tabular}
\caption{Progresivo (tasa marginal): Zapopan, Jalisco (2023).}
\label{tab:esquema-progresivo}

\end{table}


\subsection*{Tasa única}
Una sola tasa (porcentaje o al millar) aplicada uniformemente al valor catastral de todos los predios, sin categorías ni rangos.
% tasa_unica — Tasa única: Santa Maria Huatulco, Oaxaca (2014).
\begin{table}[htbp]
\centering
\small
\begin{tabular}{lrll}
\toprule
Descripción & Tasa & Base de cálculo & Unidad \\
\midrule
Impuesto Predial del 0.5\% anual sobre el valor fiscal del inmueble & 0.005 & valor\_fiscal & porcentaje \\
\bottomrule
\end{tabular}
\caption{Tasa única: Santa Maria Huatulco, Oaxaca (2014).}
\label{tab:esquema-tasa_unica}

\par\smallskip\footnotesize Mínimo predial: \$2 vsm (anual).
\end{table}


\subsection*{Cuota fija (tarifa única)}
Monto fijo en pesos por predio, independiente del valor catastral. En la tesis se denomina \emph{tarifa única}. Algunos municipios la escalonan por rangos de valor (variante \emph{cuota fija escalonada}), que aquí se agrupa dentro de esta categoría.
% cuota_fija — Cuota fija (tarifa única): Teya, Yucatán (2011).
\begin{table}[htbp]
\centering
\small
\begin{tabular}{lrll}
\toprule
Descripción & Monto & Periodicidad & Unidad \\
\midrule
Predios urbanos & \$60 & anual & pesos \\
\bottomrule
\end{tabular}
\caption{Cuota fija (tarifa única): Teya, Yucatán (2011).}
\label{tab:esquema-cuota_fija}

\end{table}


\subsection*{Mixto}
Estructura híbrida por rangos que combina columnas heterogéneas (cuotas fijas y/o tasas) según el tipo de predio, sin reducirse a las variantes anteriores.
% mixto — Mixto: Tampico, Tamaulipas (2018).
\begin{table}[htbp]
\centering
\small
\begin{tabular}{rrrr}
\toprule
Rango & Lím. inf. & Lím. sup. & Col. 1 \\
\midrule
1 & \$0.01 & \$242{,}307.69 & 0.5406 \\
2 & \$242{,}307.70 & \$500{,}000 & 0.58 \\
3 & \$500{,}000.01 & \$750{,}000 & 0.65 \\
4 & \$750{,}000.01 & \$1{,}000{,}000 & 0.73 \\
5 & \$1{,}000{,}000.01 & \$1{,}250{,}000 & 0.81 \\
6 & \$1{,}250{,}000.01 & \$1{,}500{,}000 & 0.89 \\
7 & \$1{,}500{,}000.01 & \$1{,}750{,}000 & 0.98 \\
8 & \$1{,}750{,}000.01 & \$2{,}000{,}000 & 1.07 \\
9 & \$2{,}000{,}000.01 & \$2{,}250{,}000 & 1.16 \\
10 & \$2{,}250{,}000.01 & En adelante & 1.25 \\
\bottomrule
\end{tabular}
\par\smallskip\footnotesize Columnas: Col. 1 = Urbano Suburbano Construido y Rustico Con Sin Construccion.
\caption{Mixto: Tampico, Tamaulipas (2018).}
\label{tab:esquema-mixto}

\par\smallskip\footnotesize Mínimo predial: \$3 uma (anual).
\end{table}


